\documentclass[11pt,mathserif]{beamer}

\usepackage{amsmath,amssymb,amsfonts}
\usepackage{graphicx}
\newcommand{\tiret}{\rule[0.6ex]{1.3ex}{0.22ex}}

% !! pour le francais
\usepackage[latin1]{inputenc}
\usepackage[english]{babel}
\usepackage{listings}
\usepackage{relsize}
\usepackage{color}
\usepackage{fancyvrb}
\usepackage{punk}

\newcommand{\fleche}{\alert{$\pmb{\longrightarrow}$}~~}
\newcommand{\scipy}{\textbf{Scipy}{} }
\newcommand{\numpy}{\textbf{Numpy}{} }
%\newcommand{\julia}{{\ECFAugie Julia}{} }
%\newcommand{\matlab}{{\ECFTeenSpirit Matlab}{} }
\newcommand{\julia}{{ Julia}{} }
\newcommand{\matlab}{{ Matlab}{} }
\newcommand{\octave}{\textbf{Octave}{} }
\newcommand{\scilab}{\textbf{Scilab}{} }
\newcommand{\python}{\textbf{Python}{} }



%%%%% for smiley
\def\mysmile#1{{\ooalign{\hfil\lower.06ex % a smiley face
 \hbox{$\scriptscriptstyle#1$}\hfil\crcr
 \hfil\lower.7ex\hbox{\"{}}\hfil\crcr
 \mathhexbox20D}}}


\newcommand{\bad}{\bf \textcolor{red}{\mysmile \frown}}
\newcommand{\neutral}{\bf \textcolor{blue}{\mysmile \minus}}
\newcommand{\good}{\bf \textcolor{green}{\mysmile \smile}}

\definecolor{dkgreen}{rgb}{0,0.6,0}
\definecolor{gray}{rgb}{0.5,0.5,0.5}
\definecolor{mauve}{rgb}{0.58,0,0.82}
\lstdefinelanguage{julia}
 {
 morekeywords=\color{blue}]{if,end,else, for, in, begin, type, function, return, Float64, Int64, <:},
 %morekeywords=[\color{yellow}]{->,=}
 sensitive=true,
 morecomment=[l]{\#},%
 morestring=[b]"
 }

\lstset{ %
  language=julia,
  framerule=0pt,
  basicstyle=\relsize{-2}\ttfamily,  % the size of the fonts that are used for the code
  %backgroundcolor=\color{black!10},  % choose the background color. You must add \usepackage{color}
  showspaces=false,               % show spaces adding particular underscores
  showstringspaces=false,         % underline spaces within strings
  showtabs=false,                 % show tabs within strings adding particular underscores
  %frame=single,                   % adds a frame around the code
  rulecolor=\color{black},        % if not set, the frame-color may be changed on line-breaks within not-black text (e.g. commens (green here))
  breakatwhitespace=false,        % sets if automatic breaks should only happen at whitespace
  keywordstyle=\color{blue},      % keyword style
  commentstyle=\color{dkgreen},   % comment style
  stringstyle=\color{mauve}  
}

%% ====== my beamer ==
\mode<presentation> {
\usetheme{default}    % sobre
\useinnertheme[shadow]{rounded}  % les numeros
}
\usefonttheme{structurebold}

\begin{document}

%****************************************************************
% Page de presentation 
%**************************************************************
\begin{frame}
\begin{center}
{\Large Introduction to the \julia language }
\end{center}
\vspace{1cm}
\begin{center}
\includegraphics[width=0.4\linewidth]{figures/julia.png}
\end{center}
\vspace{1cm}
\begin{center}
{\large Marc Fuentes - SED Bordeaux \\}
\end{center}
\end{frame}

%****************************************************************
% Plan 
%**************************************************************

\begin{frame}
\frametitle{Outline}
\begin{enumerate}[<+->]
\item motivations
\item \julia as a numerical language 
\item types and methods 
\item about performance
\end{enumerate}
\end{frame}


%****************************************************************
% Introduction 
%**************************************************************

\begin{frame}
\frametitle{Why yet another \matlab-like language?}
\begin{itemize}[<+->]
\item {\em scripting languages} such as \matlab, \scipy, \octave etc... are
efficient to prototype algorithms (more than C++/Fortran)
\item \matlab is not free neither open source. It still remains the {\em lingua franca}
in numerical algorithms. 
\item \octave is very slow (but highly compatible with \matlab)
\item \python along \numpy and \scipy is a beautiful language, but a little bit slow and do not
support any parallelism as a built-in feature. 
\item \textbf{Pypy} which is a nice JIT compiler of \python, but does not support neither \numpy neither \scipy
\item \textbf{R} is well suited for statistics, but suffer from old-syntax troubles
\end{itemize}
\pause
\fleche Why do not try a new language for numerical computation?
\end{frame}


%****************************************************************
% A language for numerical computation 
%**************************************************************

\begin{frame}[fragile]
\frametitle{A language for numerical computation}
\begin{itemize}[<+->]
 \item \julia's syntax is very similar to langages as \matlab, \python or \scilab, so
  switching to \julia is fast
 \item do not require {\em vectorized} code to run fast(JIT compiler)
 \item it uses references (also for function arguments)
 \item indices start to 1 and finish to \texttt{end}
 \item use brackets \texttt{[,]} for indexing
 \item it supports broadcasting
 \item support 1D arrays
\end{itemize}
\pause
\fleche let us have a look to some examples
\end{frame}

%****************************************************************
% Functional aspects
%**************************************************************

\begin{frame}[fragile]
\frametitle{Functional aspects}
\begin{itemize}[<+->]
   \item support anonymous functions : \texttt{ (x -> x*x)(2) }
   \item support \texttt{map}, \texttt{reduce}, \texttt{filter} functions
   \item functions support variadic arguments (using tuples)
   \item comprehension lists
   \item functions are not supposed to modify their arguments, otherwise
   they follow the \texttt{!} convention like \texttt{sort!}
\end{itemize}
\end{frame}

%****************************************************************
% Parallelism
%**************************************************************
\begin{frame}[fragile]
\frametitle{Parallelism}
\pause
\begin{itemize}[<+->]
   \item \julia has a built-in support for a distributed memory parallelism
   \item one-sided message passing routines 
       \begin{itemize}
            \item \texttt{remotecall} to launch a computation on a given process 
            \item \texttt{fetch} to retrieve informations 
       \end{itemize}     

   \item  high level routines
       \begin{itemize}
            \item \texttt{@parallel} reduction for lightweight iterations
            \item \texttt{pmap} for heavy iterations
       \end{itemize}     
   \item support for distributed arrays in the standard library 
\end{itemize}
\end{frame}

%****************************************************************
% external libraries calls
%**********²****************************************************
\begin{frame}[fragile]
\frametitle{external libraries calls}
\pause
\begin{itemize}[<+->]
    \item sometimes you need to call a C/Fortran code
    \item "no boilerplate" philosophy : do not require { Mexfiles}, Swig or other wrapping system
    \item the code must be in a {\bf shared} library 
    \item the syntax is the following  
\end{itemize}
\pause
\begin{lstlisting}
ccall(:function, "lib"), return_type, (type_1,...,type_n), args)
\end{lstlisting}
\end{frame}

%****************************************************************
% Types  
%**********²****************************************************

\begin{frame}[fragile]
\frametitle{Types}
\pause
\begin{itemize}[<+->]
    \item There is a graph type in Julia reflecting the hierarchy of types
\begin{lstlisting}
  None <: Int64 <: Number <: Real <: Any
\end{lstlisting}
    \item support both abstract and concrete types
    \item user can annotate the code with operator \texttt{::} "is an instance of"
    \item \julia supports 
    \begin{itemize}
        \item composite types 
        \item union types  
        \item tuple types  
        \item parametric types  
        \item singleton types
        \item type aliases
    \end{itemize}  
\end{itemize}
\end{frame}

%****************************************************************
% Multiple dispatch 
%**********²****************************************************

\begin{frame}[fragile]
\frametitle{Multiple dispatch}
\begin{itemize}
   \pause
   \item main idea : define piecewisely methods or functions depending on their arguments types
   \pause
   \item let us define \texttt{f}
   \begin{lstlisting}
   f(x::Float64,y::Float64) = 2x + y
   f(x::Int,y::Int) = 2x + y
   f(2.,3.) # returns 7.0
   f(2,3) # returns 7.0
   f(2, 3.) # throw an ERROR: no method f(Int64,Float64)
   \end{lstlisting}
   \pause
   but if we define \texttt{g}
   \begin{lstlisting}
   g(x::Number, y::Number) = 2x + y
   g(2.0, 3) # now returns 7.0
   \end{lstlisting}
   \pause
   \item no automatic or magic conversions : for operators 
    arguments are promoted to a common type (user-definable) and
    use the specific implementation
   \pause
   \item supports parametric methods
   \begin{lstlisting}
   myappend{T}(v::Vector{T}, x::T) = [v..., x] 
   \end{lstlisting}
\end{itemize}
\end{frame}

%****************************************************************
% Performance 
%**********²****************************************************

\begin{frame}[fragile]
\frametitle{Performance}
\begin{itemize}[<+->]
  \item to prove that \julia is fast language, we did some tests
  \item benchmarks sources taken from the Julia site and modified
  \item all times are in milliseconds
  \item on a Z800 (8 threads, 24G of memory), 
\end{itemize}
\pause
\begin{figure}[H]
\begin{tabular}{|c|c|c|c|c|c|c|}
\hline
appli & fib & mandel & quicksort & pisum & randstat & randmul \\
\hline
\matlab & 191 & 22 & 28 & 57 & 97 & 69 \\
Octave & 924 & 310 & 1138 & 21159 & 484 & 109 \\
Python &  4  & 7  & 14 & 1107 & 253 & 101 \\
Pypy   & 8  &  3(faux) & 13 & 44 & xxx &xxx \\ 
Julia  & 0.09 & 0.28 & 0.57 & 45 & 34 & 49 \\
Fortran & 0.08 & $\leqslant 10^{-6}$ & 0.62 & 44 & 16 & 275(16) \\
\hline
\end{tabular}
\end{figure}
\end{frame}
%**************************************************************
%  Conclusion     
%**************************************************************
\begin{frame}[fragile]{Conclusion}
\begin{itemize}[<+->] 
 \item pros : 
   \begin{itemize}
        \item  a true language, with lots of powerful features,  
        \item \julia is fast! 
   \end{itemize}
 \item cons : 
   \begin{itemize}
    \item poor graphics support (only 2D with additional package),
    \item no support for shared-memory parallelism 
    \item small community 
   \end{itemize}
 \item some non presented points : 
     \begin{itemize}
      \item meta-programming aspects : macros 
      \item reflection 
      \item packaging system based on Git
     \end{itemize} 
 \item more info at \texttt{http://julialang.org/}
\end{itemize} 
\end{frame}
\end{document}
